% Hafta 4 — Biçim dizisi açığı
% Derleme: py -3.12 tools/latex/derle.py docs/week-4/assets/h04-07-bicim-dizisi.tex
\documentclass[tikz,border=8pt]{standalone}
\usepackage{cen429-cizim}
\begin{document}
\begin{tikzpicture}[node distance=9mm and 12mm]
  \node[baslik] (b) at (0,0) {Biçim dizisi açığı};
  \node[altbaslik] at (b.south west) {tek fark: biçim dizgesini kim belirliyor};

  \node[kotukutu=62mm, below=16mm of b.south west, anchor=north west] (sol)
    {\kb{printf(girdi)}\\[2mm]
     girdi BİÇİM DİZGESİ sayılır\\[3mm]
     \%x $\rightarrow$ yığını okur\\
     \%n $\rightarrow$ belleğe YAZAR\\[2mm]
     $\rightarrow$ bilgi sızıntısı, kod akışı};
  \node[iyikutu=62mm, right=22mm of sol] (sag)
    {\kb{printf(``\%s'', girdi)}\\[2mm]
     biçim dizgesi SABİT\\[3mm]
     girdi yalnız VERİ\\[3mm]
     güvenli};

  \node[fit=(sol)(sag), inner sep=0] (ikifit) {};
  \coordinate (orta) at ($(sol.east)!0.5!(sag.west)$);
  \draw[cizgi, line width=1pt] (ikifit.north -| orta) -- (ikifit.south -| orta);

  \node[dip, text width=170mm, below=10mm of ikifit.south, anchor=north]
    {En tehlikelisi \%n: yazdığı bayt sayısını belleğe yazar $\rightarrow$ seçili adrese yazma.};
\end{tikzpicture}
\end{document}
